%!TEX TS-program = xelatex
\documentclass[a4paper,10pt,twocolumn,oneside]{article}
\setlength{\columnsep}{10pt}  
\setlength{\columnseprule}{0pt}      

\usepackage{amsthm}	
\usepackage{amssymb}
%\usepackage[margin=2cm]{geometry}
\usepackage{fontspec}
\usepackage{color}
\usepackage[x11names]{xcolor}
\usepackage{listings}	
\usepackage[Glenn]{fncychap}
\usepackage{fancyhdr}
\usepackage{graphicx}
\usepackage[export]{adjustbox}								%Graphic
\usepackage{enumerate}
\usepackage{changepage}
\usepackage{titlesec}
\usepackage{amsmath}
\usepackage{codebook}
\usepackage{verbatim}
\usepackage[CheckSingle, CJKmath]{xeCJK}
% \usepackage{CJKulem}

%\usepackage[T1]{fontenc}
\usepackage{amsmath, courier, listings, fancyhdr, graphicx}
\topmargin=0pt
\headsep=5pt
\textheight=780pt
\footskip=0pt
\voffset=-40pt
\textwidth=545pt
\marginparsep=0pt
\marginparwidth=0pt
\marginparpush=0pt
\oddsidemargin=0pt
\evensidemargin=0pt
\hoffset=-42pt

%%%%%%%%%%%%%%%%%%%%%%%%%%%%%

\setmainfont{Consolas}	
\setmonofont{Monaco}
%\setCJKmainfont{Source Han Sans}
% \setCJKmainfont{Consolas}
%\setmainfont{sourcecodepro}
\XeTeXlinebreaklocale "zh"
\XeTeXlinebreakskip = 0pt plus 1pt
\setcounter{secnumdepth}{3}

%%%%%%%%%%%%%%%%%%%%%%%%%%%%%
\makeatletter
\lst@CCPutMacro\lst@ProcessOther {"2D}{\lst@ttfamily{-{}}{-{}}}
\@empty\z@\@empty
\makeatother
\lstset{											% Code
language=C++,										% the language of the code
basicstyle=\ttfamily, 						% the size of the fonts that are used for the code
%numbers=left,										% where to put the line-numbers
numberstyle=\footnotesize,						% the size of the fonts that are used for the line-numbers
stepnumber=1,										% the step between two line-numbers. If it's 1, each line  will be numbered
numbersep=5pt,										% how far the line-numbers are from the code
backgroundcolor=\color{white},					% choose the background color. You must add \usepackage{color}
showspaces=false,									% show spaces adding particular underscores
showstringspaces=false,							% underline spaces within strings
showtabs=false,									% show tabs within strings adding particular underscores
frame=false,											% adds a frame around the code
tabsize=2,											% sets default tabsize to 2 spaces
captionpos=b,										% sets the caption-position to bottom
breaklines=true,									% sets automatic line breaking
breakatwhitespace=false,							% sets if automatic breaks should only happen at whitespace
escapeinside={\%*}{*)},							% if you want to add a comment within your code
morekeywords={*},									% if you want to add more keywords to the set
keywordstyle=\bfseries\color{Blue1},
commentstyle=\itshape\color{Red4},
stringstyle=\itshape\color{Green4},
literate={\ \ }{{\ }}1
}

\university{1}{Saratov State University}{SSU.png}
\team{ZAVRS}{Eremenko, Utaliev, Yanchenko}
\contest{Contest name}{Contest Date}

%%%%%%%%%%%%%%%%%%%%%%%%%%%%%

\begin{document}
\maketeampage

\pagestyle{fancy}
\fancyfoot{}
%\fancyfoot[R]{\includegraphics[width=12pt]{SSU.png}}
\fancyhead[L]{Saratov State University: Eremenko, Utaliev, Yanchenko}

\fancyhead[R]{\thepage}
\renewcommand{\headrulewidth}{0.4pt}
\renewcommand{\contentsname}{Contents} 
\scriptsize
\begingroup
\let\clearpage\relax
% \tableofcontents
\endgroup
\tableofcontents
%%%%%%%%%%%%%%%%%%%%%%%%%%%%%
% \newpage

% ВАЖНО!
% Для полуфинала максимальное количество страниц <= 25. (на момент 24 года).

\section{Basic}
\code{Pragma optimization}{basic/pragma.cpp}
\code{Hash Function}{basic/hash.cpp}

\section{Math}
\code{Integration}{math/integration.cpp}
\codewithtext{Simulated Annealing}{math/annealing.cpp}
{Для нахождения глобального минимума/максимума функции.} %удалим, если что
\code{Gauss}{math/gauss.cpp}
\code{Iteration}{math/iteration.cpp} %удалим, если что
\codewithtext{Tridiagonal Matrix Algorithm}{math/tridiag.cpp}{$a_{i}x_{i-1}+b_{i}x_{i}+c_{i}x_{i+1}=d_{i}$} %удалим, если что


\section{Number Theory}
\code{FFT}{nt/fft.cpp} %Нужно решить, оставить код Миши или накатить свой
\code{NTT}{nt/ntt.cpp} %Нужно решить, оставить код Миши или накатить свой
\codewithtext{FWHT}{nt/fwht.cpp}
{$c_i = \sum_{j,k:\,j \oplus k =i } a_j \cdot b_k$} %Код не проверен
\code{Extended Euclid}{nt/exgcd.cpp}
\tex{Burnside Lemma}{nt/burnside.tex}
\code{Phi and Mobius}{nt/phimu.cpp}
\codewithtext{Pollard Rho}{nt/pollard.cpp}
{Факторизация числа с малыми множителями в разложении. $O(N^{1/4})$}
\code{Miller Rabin}{nt/miller.cpp}
\code{Primitive Root}{nt/primitive.cpp}%пусть болтается
\codewithtext{Discrete Root}{nt/discrete\_root.cpp}{Найти все $x$: $x^k \equiv a\; (\text{mod}\;n)$}  %пусть болтается
\codewithtext{Discrete log}{nt/discreteLog.cpp}{$a^x \equiv  b\; (\text{mod}\;m)$}
\codewithtext{Diophantine Equations}{nt/diophantine.cpp}{solutions to $ax + by = c$ where $x \in [xlow, xhigh]$ and $y \in [ylow, yhigh]$; returns \{cnt, leftsol, rightsol, gcd of a and b\}}
\tex{Stern-Brocot Tree}{nt/sternBrocot.tex}
\codewithtext{Lattice Points Below Line}{nt/lattice\_points.cpp}{Number of integer points $(x;y)$ such for $0 \leq x < n$ and $0 < y \leq \lfloor k x+b\rfloor$}
%добавить структуру для рациональных чисел. Нужно ли?? -нет конечно

\section{Graph Theory}
\code{Biconnected Components}{graph/bicon.cpp} %Нужно ли block-cut tree?
\code{2-SAT}{graph/twosat.cpp}
\tex{Kirchoffs Theorem}{graph/kirghoff.tex}
\tex{Tuttes Theorem}{graph/tutte.tex}
\tex{Matching Duals}{graph/matching\_dual.tex}
\code{LCA}{graph/lca.cpp}
\code{Kuhn}{graph/kuhn.cpp}
\code{1-K bfs}{graph/1kBFS.cpp}
\code{Floyd-Warshall}{graph/floydWarshall.cpp}
\code{Ford-Bellman with negative cycle}{graph/FordBellman.cpp}
\codewithtext{Eulerian path}{graph/eulerianPath.cpp}
{
Эйлеров цикл существует тогда и только тогда, когда степени всех вершин чётны. Эйлеров путь существует тогда и только тогда, когда количество вершин с нечётными степенями равно двум (или нулю). Граф должен быть достаточно связным (если удалить из него все изолированные вершины, то должен получиться связный граф).

Если нечетных вершин ровно две, то соединим их ребром, построим эйлеров цикл, затем удалим это ребро из цикла. Если правильно сдвинуть этот цикл, получим эйлеров путь.
}
% Добавить Борувко

\section{Flows}
\code{Dinics Algorithm}{flows/dinics.cpp}
\code{MCMF}{flows/mcmf.cpp}
%следующее оставляем как черный ящик, если что, удалим
\codewithtext{MCMF with Potentials}{flows/mcmf\_potentials.cpp}{call init\_dag() or init\_fb() (depending on if your graph is dag or not) before running calc()} 
\tex{L-R Flow}{flows/lrflow.tex}
\code{Stoer Wagner Algorithm}{flows/stoer-wagner.cpp} %нужно написать разъяснения, ограничение и время
\code{Hungarian Algorithm}{flows/hungarian.cpp} %нужно написать че делает
\code{Blossom Algorithm}{flows/blossom.cpp} %нужно написать че делает
%нужно добавить задачу про проекты и инструменты


\section{Data Structures}
\code{Centroid Decomposition}{ds/centroid.cpp} %нужно разобраться
\code{HLD}{ds/hld.cpp}
\code{Explicit Treap}{ds/treap.cpp}
\code{Sparce Table}{ds/sparceTable.cpp}
\code{Persistent Segment Tree}{ds/persDO.cpp}
%добавить персистентный декартач
\code{Ordered Set}{ds/pbds.cpp}


\section{Strings}
\tex{String Matching with FFT}{strings/fft\_match.tex}
\codewithtext{Suffix Array}{strings/suffix\_array.cpp}{Don't forget to modify the string to avoid cyclic comparisons if needed}
\code{Prefix- and z-funtion}{strings/prefix-z.cpp}
\code{Manachers Algorithm}{strings/manacher.cpp}
\code{Prefix Automaton}{strings/prefix\_automaton.cpp}
\code{Aho-Corasick + Trie}{strings/aho.cpp}
\code{Hash String}{strings/hash.cpp}


\section{DP}
\code{Dynamic CHT}{dp/cht.cpp} %пояснить че делает
\code{Li Chao Tree}{dp/lichao.cpp} %надо бы разобраться
\code{D\&C Optimization}{dp/dnc.cpp} %надо написать условия
%нужно переписать условие
\codewithtext{Knuth Optimization}{dp/knuth.cpp}{Достаточное условие: 1. $C_{ac}+C_{bd}\le C_{ad}+C_{bc}$, 2. $C_{bc}\le C_{ad}$, при всех $a\le b\le c\le d$.}
%это надо переписать
\codewithtext{Slope Trick}{dp/slope.cpp}{You are given an array, each operation you are allowed to increase or decrease an element's value by $1$. Find the minimum number of operations to make the array strictly increasing.}
%добавить лямбда оптимизацию
\codewithtext{SOS DP}{dp/sosdp.cpp}
{or-свертка:$c_i = \sum_{j,k:\,j | k = i} a_j \cdot b_k$. Нужно взять сумму по подмножествам пос-тей $a$ и $b$, перемножить поэлементно, а потом обратить суммы по подмножествам для этого произведения.

and-свертка: $c_i = \sum_{j,k:\,j \& k = i} a_j \cdot b_k$. Нужно взять сумму по \textbf{над}множествам пос-тей $a$ и $b$, перемножить поэлементно, а потом обратить суммы по \textbf{над}множествам для этого произведения.}
\code{Submasks}{dp/submasks.cpp}
%добавить нвп бинпоиском


\section{Geometry}
\code{Base}{geometry/base.cpp}
\code{Intersections}{geometry/intersections.cpp}
\code{Angles}{geometry/angles.cpp}
\code{Areas}{geometry/areas.cpp}
\code{Distances}{geometry/distances.cpp}
\code{Circle Circle Area}{geometry/circle\_circle\_area.cpp}
\code{Circle Line Intersection}{geometry/circle\_line.cpp}
\code{Circle Circle Intersection}{geometry/circle\_circle.cpp}
\code{Tangent Lines of Two Circles}{geometry/tangents\_circles.cpp}
\code{Convex Hull}{geometry/convex.cpp}
\codewithtext{Minkowski Sum}{geometry/minkowski.cpp}{Рассмотрим два множества \( A \) и \( B \) точек на плоскости. Сумма Минковского \( A + B \) определяется как \( \{a + b \mid a \in A, b \in B\} \). Будем рассматривать случай, когда \( A \) и \( B \) состоят из выпуклых многоугольников \( P \) и \( Q \) с их внутренностями. Cумма выпуклых многоугольников \( P \) и \( Q \) является выпуклым многоугольником с не более чем \( |P| + |Q| \) вершинами.}
\code{Point in Convex Polygon}{geometry/point\_in\_convex.cpp}
\code{SVG}{geometry/svg.cpp}


\section{Miscellaneous}
\codewithtext{Josephus Problem}{misc/josephus.cpp}{Группа из $n$ людей стоит в круге, и каждый $k$-й человек по счёту исключается из круга до тех пор, пока не останется один человек. Задача заключается в определении позиции этого последнего оставшегося человека.} 
\codewithtext{Knight Moves in Infinity Grid}{misc/knight.cpp}{Задача заключается в определении минимального количества ходов, которые необходимы коню для достижения заданной позиции на бесконечной шахматной доске, начиная с определённой исходной позиции.}

\setlength{\parindent}{0cm}


\section{Other}
%можно добавить факты про псп
\tex{Комбинаторика}{other/combinatorics.tex}
%нужно добавить сумму делителей
\tex{Теория чисел}{other/number-theory.tex}
\tex{Геометрия}{other/geometry.tex}
\tex{Графы}{other/graphs.tex}
%нужно переписать интерполяционный многочлен ньютона
\tex{Формулы}{other/other.tex}
\tex{Полезные числа}{other/table.tex}

\end{document}